% Part: computability
% Chapter: recursive-functions
% Section: bounded-minimization

\documentclass[../../../include/open-logic-section]{subfiles}

\begin{document}

\olfileid{cmp}{rec}{bmi}
\olsection{সীমাবদ্ধ ন্যূনীকরণ}

\begin{explain}
কোনো ধর্ম বা সম্বন্ধ~$P$ পূরণকারী ক্ষুদ্রতম সংখ্যা হিসেবে একটি অপেক্ষককে
সংজ্ঞায়িত করা প্রায়ই উপযোগী। $P$ নির্ণেয় হলে, $P$ পূরণকারী ক্ষুদ্রতম
সংখ্যাটি না পাওয়া পর্যন্ত সম্ভাব্য সব সংখ্যা $0$, $1$, $2$, \dots,
পরীক্ষা করে আমরা সহজেই এই অপেক্ষকটি গণনা করতে পারি। এই ধরনের সীমাহীন
অনুসন্ধান আমাদের আদিম পুনরাবৃত্ত অপেক্ষকের পরিসরের বাইরে নিয়ে যায়। কিন্তু
যদি কেবল \emph{স্বাধীনভাবে প্রদত্ত কোনো সীমার চেয়ে ছোট} ক্ষুদ্রতম
সংখ্যাটিতে আমাদের আগ্রহ থাকে, তবে আমরা আদিম পুনরাবৃত্তির পরিসরেই থাকি।
অন্যভাবে এবং একটু সাধারণভাবে বলা যাক, আমাদের একটি আদিম পুনরাবৃত্ত সম্বন্ধ
থাকুক~$R(x,z)$। যে অপেক্ষকটি $x$ ও~$y$-কে এমন ক্ষুদ্রতম $z < y$-তে
পাঠায় যার জন্য $R(x, z)$ সত্য, সেটি বিবেচনা করুন। $R(x, 0)$,
$R(x, 1)$, \dots, $R(x, y-1)$ সত্য কি না পরীক্ষা করে এটিকেও গণনা করা
যায়। কিন্তু এটি আদিম পুনরাবৃত্ত কেন?
\end{explain}

\begin{prop}
$R(\vec x, z)$ আদিম পুনরাবৃত্ত হলে $m_R(\vec{x}, y)$ অপেক্ষকটিও আদিম
পুনরাবৃত্ত; এমন সংখ্যা থাকলে এটি $y$-এর চেয়ে ছোট এমন ক্ষুদ্রতম~$z$ ফেরত
দেয় যার জন্য $R(\vec x, z)$ সত্য, আর এমন সংখ্যা না থাকলে $y$ ফেরত দেয়।
অপেক্ষক~$m_R$-কে আমরা লিখব
\[
\bmin{z < y}{R(\vec{x}, z)},
\]
\end{prop}

\begin{proof}
$R(\vec{x}, z)$ সত্য হয় এমন কোনো~$z < 0$ থাকতে পারে না, কারণ আদৌ
কোনো $z < 0$ নেই। অতএব $m_R(\vec x, 0) = 0$।

সীমাটি $y + 1$ রূপের হলে তিনটি ক্ষেত্র হয়:
\begin{enumerate}
\item এমন একটি $z < y$ আছে যার জন্য $R(\vec{x}, z)$ সত্য; সেক্ষেত্রে
$m_R(\vec{x}, y+1) = m_R(\vec{x}, y)$।
\item এমন কোনো~$z<y$ নেই, কিন্তু $R(\vec{x}, y)$ সত্য; সেক্ষেত্রে
$m_R(\vec{x}, y+1) = y$।
\item এমন কোনো $z < y+1$ নেই যার জন্য $R(\vec{x}, z)$ সত্য; সেক্ষেত্রে
$m_R(\vec{x}, y+1) = y+1$।
\end{enumerate}
% BN-SRC-140 TARGET-CORRECTION-BEGIN
\par\smallskip\noindent\textnormal{\small\textbf{উৎস-সংশোধন (BN-SRC-140):} হিমায়িত ইংরেজি পাঠে তৃতীয় ক্ষেত্রের বাম পাশে আগের সব ক্ষেত্রের আর্গুমেন্ট-ভেক্টরের বদলে ভুল করে z-ভেক্টর লেখা হয়েছে। বাংলা সূত্রে একই x-ভেক্টর পুনরুদ্ধার করা হয়েছে।} % BN-SRC-140 TARGET-CORRECTION-END
তাই আদিম পুনরাবৃত্তির সাহায্যে $m_R(\vec x, 0)$-কে নিচের মতো
সংজ্ঞায়িত করা যায়:
\begin{align*}
m_R(\vec{x}, 0) & = 0\\
m_R(\vec{x}, y+1) & =
\begin{cases}
m_R(\vec{x}, y) & \text{যদি $m_R(\vec x, y) \neq y$}\\
y & \text{যদি $m_R(\vec x, y) = y$ এবং $R(\vec{x}, y)$}\\
y+1 & \text{অন্যথায়।}
\end{cases}
\end{align*}
লক্ষ করুন, এমন একটি $z<y$ আছে যার জন্য $R(\vec x, z)$ সত্য, যদি এবং
কেবল যদি $m_R(\vec x,
y) \neq y$।
\end{proof}

\begin{prob}
ধরা যাক, $R(\vec x, z)$ আদিম পুনরাবৃত্ত। আদিম পুনরাবৃত্তির সাহায্যে
~$\Char{R}$ থেকে $m'_R(\vec{x}, y)$ অপেক্ষকটি সংজ্ঞায়িত করুন, যা
এমন সংখ্যা থাকলে $y$-এর চেয়ে ছোট এমন ক্ষুদ্রতম~$z$ ফেরত দেয় যার জন্য
$R(\vec{x}, z)$ সত্য, আর এমন সংখ্যা না থাকলে $0$ ফেরত দেয়।
\end{prob}

\end{document}
